% !TeX program = lualatex
% !TeX encoding = utf8
% !TeX spellcheck = uk_UA
% !TeX root = ../main.tex

%% --------------------------------------------------------
%% Chapter file: Nature.tex (розділ 0: огляд, без задач)
%% --------------------------------------------------------

\graphicspath{{\currfilebase/Pictures}}

\setcounter{chapter}{-1}
\chapter{Плазма довкола нас}\label{sec:plasma-nature}\hypertarget{plasma-nature}{}

Почнімо з~того, що просто подивимося на плазму: числа й~формули чекатимуть у~наступних розділах. Спершу домовимося, що саме називаємо плазмою, а~далі здійснимо екскурсію від космосу до лабораторії й~техніки.
За оцінками, понад 99\,\% звичайної (барионної) видимої речовини у~Всесвіті перебуває в~стані плазми. У~цьому розділі нічого не виводимо: лише розглядаємо знімки й~вчимося бачити на них фізику,
яку з'ясуємо в~наступних розділах. Частину об'єктів (іоносферу, корону, сонячний вітер, газовий розряд, токамак) знайдемо в~табл.~\ref{tab:plasma-types}
і~на діаграмі рис.~\ref{fig:plasma-types} у~§\,\ref{sec:nature-common}.
Якщо якесь слово чи число здається незрозумілим, не зупиняйтеся: посилання вказує розділ, де це пояснено.

\section{Що таке плазма?}\label{sec:what-is-plasma}

\emphx{Плазма}\index{Плазма|textbf} (від грец. \emphx{plasma}~--- виліплене, оформлене)~--- це частково або
повністю іонізований газ, у~якому густини позитивних і~негативних зарядів практично
однакові, а~заряджені частинки поводяться \emphx{колективно}\index{Колективність|textbf}. Обидва виділені слова мають
точний зміст, який ми з'ясуємо в~розд.~\ref{sec:plasma-basics} (квазінейтральність~--- у~§\,\ref{sec:quasineutrality},
колективність~--- у~§\,\ref{sec:debye-sphere}). Термін \enquote{плазма} запропонував
І.~Ленгмюр наприкінці 1920-х років; 1929~р. Л.~Тонкс і~Ленгмюр описали власні коливання
електронного газу (§\,\ref{sec:plasma-oscillations})~\cite{chen2016introduction}.

Ступінь іонізації характеризують величиною
\begin{equation}\label{eq:ionization-degree}
	\alpha = \frac{n_i}{n_i + n_n},
\end{equation}
де $n_i$ і~$n_n$~--- концентрації іонів та нейтральних частинок. У~плазмі можуть одночасно
бути присутніми електрони, кілька сортів іонів і~нейтральні атоми чи молекули. Залежно від
$\alpha$ розрізняють \emphx{слабо іонізовану}\index{Плазма!слабоіонізована|textbf} (частки відсотка), \emphx{помірно
іонізовану} (кілька відсотків) і~\emphx{повністю іонізовану}\index{Плазма!повністю іонізована|textbf} ($\alpha\to1$) плазму.
Газ іонізується помітно вже за температур, значно нижчих за енергію іонізації атома (для водню $13{,}6$~еВ); чому, пояснює формула Саха (§\,\ref{sec:saha}).

\emphx{Чим плазма відрізняється від нейтрального газу?} Не просто наявністю заряджених
частинок. У~нейтральному газі частинки майже вільні й~обмінюються імпульсом
у~короткочасних зіткненнях. У~плазмі кулонівська сила спадає повільно ($\propto r^{-2}$): кількість частинок у~сферичному шарі
радіуса $r$ зростає як $r^2$, тож для впорядкованого (усередненого) відгуку сумарний внесок далеких частинок не менший, ніж близьких,
і~кожна заряджена частинка \emph{одночасно} взаємодіє з~великою кількістю інших. Сила,
що діє на частинку, формується усередненим полем багатьох частинок, а~не парними
зіткненнями. Такі \enquote{колективні взаємодії} породжують екранування зарядів, плазмові
коливання, хвилі й~нестійкості. Тому плазма не зводиться до сукупності незалежних
частинок: її макроскопічні властивості визначаються спільною електромагнітною взаємодією.
Ступінь колективності задає число $N_D$ частинок у~сфері радіуса $\lambda_D$
(§\,\ref{sec:debye-sphere}): поле заряду на відстанях понад дебаївський радіус гасить не найближчий
сусід, а~спільний внесок цих $N_D$ частинок (§\,\ref{sec:debye-screening}).

Друга особливість~--- тісний зв'язок механіки з~електродинамікою. Плазма не лише рухається
в~зовнішньому електромагнітному полі, а~й~сама створює макроскопічні поля та струми, які
діють назад на її рух.

\begin{Theory}[Температура в~енергетичних одиницях]
Температуру $T$ вимірюємо в~енергетичних одиницях ($k_{\mathrm B}=1$, тож $k_{\mathrm B}T$ пишемо просто
як $T$); практично $T$ дають в~еВ, $1$~еВ $\approx11\,600$~К. Температура~--- це параметр
максвеллівського розподілу $f\propto\exp(-mv^2/2T)$; середня кінетична енергія поступального руху
$\langle mv^2/2\rangle=\tfrac32T$. Поки не сказано інакше, мова йде саме про такий розподіл.
Для оцінок зручні числа (СГС): $e=4{,}8\times10^{-10}$~СГС, $e^2\approx1{,}44\times10^{-7}$~еВ$\cdot$см,
$1$~СГС напруженості поля $=300$~В/см.
\end{Theory}

З~термодинаміки відомо, що стан газу задають $p$, $V$ і~$T$, але незалежних серед них лише два: їх пов'язує рівняння стану, для ідеального газу $pV=NT$ (при $k_{\mathrm B}=1$). Зручно взяти інтенсивні величини, що не залежать від кількості речовини: $n=N/V$ і~$T$; тоді тиск $p=nT$ визначається ними. Так само й~плазму (заданого складу) характеризують двома параметрами, $n$ й~$T$ (для суміші сортів $p=\sum_s n_sT_s$). До того ж від $n$ і~$T$ залежить уся специфічна для плазми фізика: ступінь іонізації (§\,\ref{sec:saha}), дебаївська довжина (§\,\ref{sec:debye-screening}), число частинок у~дебаївській сфері (§\,\ref{sec:debye-sphere}), плазмова частота (§\,\ref{sec:plasma-oscillations}). Тому плазму зображають точкою на діаграмі $n$--$T$ (рис.~\ref{fig:plasma-types}); докладніше про температуру~--- у~§\,\ref{sec:two-temperatures}.

Плазма виникає різними способами: термічною іонізацією (нагрівання), електричним розрядом,
фотоіонізацією, ударною іонізацією, бомбардуванням частинками. Залежно від температури її
поділяють на \emphx{низькотемпературну}\index{Плазма!низькотемпературна|textbf} ($T\lesssim10^5$~К, тобто $\lesssim10$~еВ) і~\emphx{високотемпературну}\index{Плазма!високотемпературна|textbf}
($T\gtrsim10^6$~К, тобто $\gtrsim100$~еВ; у~термоядерних установках (токамаках~--- тороїдальних магнітних камерах)~--- до $10^8$~К~$\approx10$~кеВ і~вище). Область між цими межами називають проміжною. Саме високотемпературна плазма потрібна для реакцій
\emphx{керованого термоядерного синтезу}. Не варто плутати цей поділ з~назвами \enquote{холодна}
і~\enquote{гаряча} плазма, які трапляються далі в~теорії хвиль: там ідеться не про саме значення $T$,
а~про те, чи істотний тепловий рух частинок для процесу, що розглядається (це пояснено там, де такі
назви з'являються, і~в~\hyperref[sec:glossary]{Словнику термінів} (с.~\pageref{sec:glossary})).
Скільки в~природі різних плазм і~де вони лежать за $n$ і~$T$, показано в~§\,\ref{sec:nature-common}.

\section{Природна плазма}\label{sec:nature-natural}

\subsection*{Всесвіт як плазма}

Почнімо з~самого початку. Плазма~--- не екзотика, а~первинний стан видимої речовини. Приблизно перші 380~тисяч років після Великого вибуху
Всесвіт був заповнений гарячою \emph{воднево-гелієвою плазмою}: протони (разом із~ядрами гелію це приблизно вся звичайна речовина), електрони й~фотони.
Фотони безперервно розсіювалися на вільних електронах, тож Всесвіт був непрозорий, як сяюча куля. Розширення охолоджувало плазму, і~коли її температура впала
до $\sim3\times10^{3}$~К (близько $0{,}3$~еВ, у~$\sim50$ разів нижче за потенціал іонізації водню $13{,}6$~еВ), електрони об'єдналися з~ядрами в~нейтральні атоми:
настала \emphx{рекомбінація}\index{Рекомбінація космологічна|textbf}. Відтоді світло летить вільно й~дійшло до нас як
\emphx{реліктове випромінювання}\index{Реліктове випромінювання|textbf}, охолоджене розширенням до $2{,}725$~К (рис.~\ref{fig:nature-cmb}).

\begin{Estimate}[Коли Всесвіт рекомбінував]
Сьогодні концентрація атомів водню у~Всесвіті~--- $n_{H}\approx2\times10^{-7}~\mathrm{см}^{-3}$. Під час рекомбінації Всесвіт був менший у~$1+z\approx1100$ разів
за кожним виміром ($z$~--- червоне зміщення), тож $n\approx2\times10^{-7}\cdot1100^{3}\approx2{,}7\times10^{2}~\mathrm{см}^{-3}$. За формулою Саха (§\,\ref{sec:saha}) для такої концентрації
ступінь іонізації $\alpha$ (§\,\ref{sec:what-is-plasma}) досягає $1/2$ за $T\approx3{,}7\times10^{3}$~К ($0{,}32$~еВ), а~за $T=3\times10^{3}$~К вона вже лише $\alpha\sim5\times10^{-3}$.
Температура реліктового випромінювання тоді була $2{,}725\cdot1101\approx3\times10^{3}$~К, тобто збігається з~оцінкою за порядком.
\end{Estimate}
Рівноважна формула Саха дає дещо вищу температуру, ніж справжня: у~Всесвіті, що розширюється, рекомбінація не встигає за рівновагою.
Важливо, що водень рекомбінує за $T\approx0{,}3$~еВ, а~не за $13{,}6$~еВ: на один атом припадає близько $10^{9}$ фотонів, і~навіть рідкісні фотони з~\enquote{хвоста} планківського розподілу з~енергією понад $13{,}6$~еВ
встигають іонізувати атоми; це той самий виграш ентропії вільних електронів, що й~у~§\,\ref{sec:saha}.

\begin{figure}[tb]\centering
	\includegraphics[width=0.85\linewidth]{cmb-wmap}
	\caption{Реліктове випромінювання: карта всього неба за даними WMAP (дев'ятирічні спостереження; NASA/WMAP Science Team, суспільне надбання).
	Кольори показують відхилення температури від середньої $2{,}725$~К в~межах $\pm2\times10^{-4}$~К (синій~--- холодніше, червоний~--- тепліше)}
	\label{fig:nature-cmb}
\end{figure}

Неоднорідності на цій карті малі, порядку $10^{-5}$ від середньої температури, але саме вони важливі: це відбитки \emph{звукових хвиль у~первинній плазмі}.
Гравітація (передусім темної матерії) стискала плазму, тиск випромінювання розпирав її, і~так вона коливалася, як струна (про коливання й~хвилі в~плазмі йтиметься в~розд.~\ref{sec:plasma-waves}).
Зі~зменшенням температури коливання \enquote{заморозилися}, а~потім з~цих згущень виросли галактики. Пізніше випромінювання перших зір знову іонізувало газ
(протягом першого мільярда років), і~міжгалактичний газ, у~якому зосереджена значна частина звичайної речовини, відтоді лишається плазмою~---
вкрай розрідженою ($n\lesssim10^{-5}~\mathrm{см}^{-3}$), з~температурою від $10^{4}$ до $10^{7}$~К.

\subsection*{Масштаби й~структури}

Природна плазма охоплює розміри від сантиметрів (канал блискавки) до світлових років (туманності) і~концентрації
від $\sim10^{-5}~\mathrm{см}^{-3}$ до $10^{26}~\mathrm{см}^{-3}$. У~розрідженій космічній плазмі зіткнень мало, а~магнітне поле істотне.
Тому на рис.~\ref{fig:nature-natural}\,а,\,б,\,в ми бачимо переважно не \enquote{газ}, а~\emph{структури, задані полем}. Об'єкти нижче впорядковано за масштабом: від центрів галактик через міжзоряне середовище й~Сонячну систему до навколоземного космосу й~атмосфери.

\begin{figure}[tb]\centering
	\begin{subfigure}[t]{0.45\linewidth}\centering
		\includegraphics[width=\linewidth]{crab-nebula}
		\subcaption{Крабоподібна туманність (Hubble; NASA, ESA, J.~Hester, A.~Loll)}
	\end{subfigure}\hfill
	\begin{subfigure}[t]{0.45\linewidth}\centering
		\includegraphics[width=\linewidth]{sun-corona}
		\subcaption{Корона Сонця, 171~\AA\ (SDO/AIA; NASA/SDO)}
	\end{subfigure}

	\medskip
	\begin{subfigure}[t]{0.45\linewidth}\centering
		\includegraphics[width=\linewidth]{aurora}
		\subcaption{Полярне сяйво з~борту МКС (NASA)}
	\end{subfigure}\hfill
	\begin{subfigure}[t]{0.45\linewidth}\centering
		\includegraphics[width=\linewidth]{lightning}
		\subcaption{Блискавка, Техас (NOAA)}
	\end{subfigure}
	\caption{Природна плазма від далекого космосу до Землі: залишок наднової~(а), корона Сонця~(б), полярне сяйво~(в), блискавка~(г). Знімки є суспільним надбанням}
	\label{fig:nature-natural}
\end{figure}

\subsection*{Дальній космос: чорні діри та джети}

Сама чорна діра не є плазмою, але плазма довкола неї~--- одна з~найекстремальніших у~природі. Речовина, що падає в~діру, не потрапляє в~неї одразу: вона закручується в~\emphx{акреційний диск}\index{Акреційний диск|textbf}
з~іонізованого газу, який нагрівається тертям до температур від мільйонів до мільярдів кельвінів. Цей газ намагнічений: магнітне поле \enquote{вморожене} в~плазму (розд.~\ref{sec:mhd}),
а~нестійкості плазми (розд.~\ref{sec:instability}) переносять момент імпульсу назовні й~дають речовині змогу падати. Саме таку картину фіксує
Event Horizon Telescope~--- мережа радіотелескопів, що працює як єдиний телескоп розміром із~Землю. Кільце на знімку 2019~р. (рис.~\ref{fig:nature-bh}\,а)~---
випромінювання гарячої плазми біля центральної діри галактики M87, а~темна пляма всередині~--- \emph{тінь} діри: світло, що летить до неї, не повертається.
Маса цієї діри близько $6{,}5\times10^{9}$ мас Сонця, тому її радіус Шварцшильда $2GM/c^{2}\approx2\times10^{13}$~м, тобто близько 130 астрономічних одиниць. Діра в~центрі нашої Галактики Sgr~A$^{*}$ у~півтори тисячі разів легша, близько 4~млн мас Сонця.

\begin{figure}[tb]\centering
	\begin{subfigure}[t]{0.45\linewidth}\centering
		\includegraphics[width=\linewidth]{bh-m87}
		\subcaption{M87$^{*}$ (EHT Collaboration, CC BY 4.0)}
	\end{subfigure}\hfill
	\begin{subfigure}[t]{0.45\linewidth}\centering
		\includegraphics[width=\linewidth]{bh-jet}
		\subcaption{Джет галактики M87 (Hubble; NASA/ESA, F.~D.~Macchetto)}
	\end{subfigure}
	\caption{Плазма навколо чорної діри: диск і~тінь M87$^{*}$~(а), джет тієї самої галактики~(б). Знімок (б)~--- суспільне надбання}
	\label{fig:nature-bh}
\end{figure}

З~околиць діри вилітають \emphx{джети}\index{Джет|textbf}~--- вузькі потоки плазми зі~швидкостями, близькими до швидкості світла, що тягнуться вздовж осі обертання на тисячі світлових років
(рис.~\ref{fig:nature-bh}\,б: джет M87 завдовжки близько 4~тисяч світлових років). Їх колімує й, імовірно, живить магнітне поле, закручене разом із~диском.
Зображення діри з~фільму \enquote{Інтерстеллар} теж обчислене за рівняннями теорії відносності, але зі~спрощеннями заради ефектності; реальні знімки менш яскраві, зате це справжня плазма.

\subsection*{Галактика: міжзоряний газ}
Ближче до нас~--- \emphx{міжзоряний газ}. Його концентрація мізерна (від $\sim10^{-3}$ до $10^{3}$ частинок в~одному кубічному сантиметрі), температура змінюється від $\sim10$~К у~холодних хмарах до $\sim10^{6}$~К у~гарячих
бульбашках, а~зіткнень настільки мало, що газ не врівноважений із~випромінюванням, а~за фронтами ударних хвиль наднових іони й~електрони довго мають різні температури. Іонізують газ ультрафіолетом зірок, ударними хвилями
наднових і~космічними променями. Поруч із~ним~--- залишки наднових, як-от Крабоподібна туманність (рис.~\ref{fig:nature-natural}\,а): наднова 1054~р., розліт газу на кілька світлових років.
Її волокна тримає магнітне поле, \enquote{вморожене} в~газ, а~блакитне світіння в~центрі~--- синхротронне випромінювання, тобто випромінювання релятивістських електронів, що крутяться в~магнітному полі туманності, яке живить пульсар (швидко обертова нейтронна зірка; розд.~\ref{sec:particle-motion}).

\subsection*{Сонячна система: Сонце й~сонячний вітер}

Сонце (рис.~\ref{fig:nature-natural}\,б)~--- найближча до нас зірка й~найдоступніша для спостережень лабораторія плазми.
У~його ядрі, де $T\sim1{,}5\times10^{7}$~К, ідуть термоядерні реакції (розд.~\ref{sec:fusion}); енергія повільно виходить назовні променистим перенесенням,
а~в~зовнішніх шарах~--- конвекцією. Рухи провідної плазми в~конвективній зоні підтримують магнітне поле Сонця: у~сонячних плямах
воно досягає $\sim0{,}4$~Тл, а~приблизно кожні 11~років змінює напрям. Поверхня (фотосфера) має $T\approx5800$~К, а~от розріджена \emphx{корона}\index{Корона Сонця|textbf}
над нею гарячіша на два порядки, $T\sim10^6$~К. Цей парадокс (звідки корона бере енергію?) остаточно не розв'язано~\cite{klimchuk2015key}; один із претендентів~--- хвилі в~плазмі (розд.~\ref{sec:mhd-waves}).
Яскраві петлі на знімку~--- це плазма, утримана в~магнітних трубках~--- пучках силових ліній~$\vect{B}$: заряджені частинки легко рухаються вздовж поля й~майже не рухаються впоперек
(розд.~\ref{sec:particle-motion}).

Корона не утримується: вона безперервно розширюється, і~потік заряджених частинок, що тече з~неї, називають
\emphx{сонячним вітром}\index{Вітер!сонячний|textbf}. Біля орбіти Землі його швидкість становить $300$--$800$~км/с, концентрація~--- кілька частинок
в~одному кубічному сантиметрі, а~температура~--- порядку $10^{5}$~К (близько 10~еВ). Вітер захоплює з~собою магнітне поле Сонця (воно \enquote{вморожене} в~плазму, розд.~\ref{sec:mhd}),
і~через обертання Сонця силові лінії закручуються в~спіраль. Область, у~якій панує сонячний вітер, називають геліосферою.

\begin{Estimate}[Скільки маси Сонце втрачає на вітер]
Потік маси крізь сферу радіуса $r=1$~а.о.\ $=1{,}5\times10^{13}$~см при $n\approx9~\mathrm{см}^{-3}$ і~$v\approx470$~км/с (протони, $m_p=1{,}67\times10^{-24}$~г):
$\dot M=4\pi r^2 n m_p v\approx2\times10^{12}$~г/с~$=2\times10^{9}$~кг/с. За рік це $6\times10^{16}$~кг, тобто $\sim3\times10^{-14}$ маси Сонця.
Вітер зі~швидкістю $400$~км/с долітає до Землі за $r/v\approx4$~доби, а~корональний викид маси, що летить зі~швидкістю близько $10^3$~км/с, за~\mbox{$\sim1{,}7$~доби}.
\end{Estimate}
Це мізерна втрата: вона порівнянна з~масою, яку Сонце щосекунди перетворює на випромінювання, але за мільярди років не змінює його масу помітно.
Зате вітер визначає все, що відбувається в~міжпланетному просторі.

\subsection*{Навколоземний космос: магнітосфера й~іоносфера}

Зустрівши магнітне поле Землі ($B\sim0{,}5$~Гс на поверхні), сонячний вітер не проходить крізь нього, а~обтікає, стискаючи його з~боку Сонця й~витягуючи у~довгий хвіст
з~нічного боку. Область, у~якій поле Землі керує рухом плазми, називають \emphx{магнітосферою}\index{Магнітосфера|textbf}. Перед нею утворюється стрибок густини й~поля~---
\emphx{ударна хвиля}\index{Ударна хвиля}: вітер надзвуковий (його швидкість $\sim400$~км/с набагато більша за швидкість магнітного звуку $\sim50$--$70$~км/с), причому хвиля така, що частинки майже не стикаються: енергію розсіюють колективні поля, а~не зіткнення
(розд.~\ref{sec:plasma-collisions}). Заряджені частинки, що потрапили в~дипольне поле Землі, відбиваються від його згущень біля полюсів і~накопичуються
в~\emphx{радіаційних поясах}\index{Радіаційні пояси|textbf} (§\,\ref{sec:magnetic-mirror}); під час магнітних бур електрони з~хвоста магнітосфери, прискорені до кількох кеВ, висипаються вздовж силових ліній в~атмосферу на висоти
$80$--$150$~км, і~виникає полярне сяйво (рис.~\ref{fig:nature-natural}\,в). Спричинені вітром магнітні бурі можуть наводити в~довгих лініях електропередач струми, що виводять з~ладу обладнання
(так, у~березні 1989~р. магнітна буря призвела до багатогодинного вимкнення електроенергії у~канадській провінції Квебек), а~також порушують зв'язок і~роботу супутників.
Прогнозування таких подій називають \emphx{космічною погодою}\index{Космічна погода|textbf}.

Нижче, на висотах від $\sim60$~км, лежить \emphx{іоносфера}\index{Іоносфера}~--- верхня атмосфера, іонізована випромінюванням Сонця. Її поділяють на шари D, E і~F; найбільша концентрація електронів
($10^{5}$--$10^{6}~\mathrm{см}^{-3}$) припадає на висоти $250$--$400$~км. Вночі, коли Сонце не іонізує газ, нижні шари майже зникають.
Для практики важливо, що така плазма відбиває радіохвилі, частота яких нижча за \emph{плазмову} $f_p$~--- власну частоту коливань електронів (§\,\ref{sec:plasma-oscillations}):
$f_p\approx9\sqrt{n}$~кГц, де $n$ у~$\mathrm{см}^{-3}$, тобто $\sim9$~МГц за $n=10^{6}~\mathrm{см}^{-3}$.
Такі хвилі не проходять, і~завдяки відбиттю від іоносфери можна приймати короткохвильові станції за тисячі кілометрів (§\,\ref{sec:ionosphere}).

\subsection*{Атмосфера Землі: блискавка}

\emphx{Блискавка} (рис.~\ref{fig:nature-natural}\,г)~--- гігантська іскра. У~грозовій хмарі більші й~важчі краплі й~льодинки заряджаються переважно негативно й~падають швидше за дрібні позитивні,
тож хмара стає електричним диполем. Коли поле десь (імовірно, біля крапель і~льодинок) стає достатнім для пробою, починається електронна лавина; з~лавин зливаються провідні канали~---
\emph{лідери}, що ступінчасто, стрибками в~десятки метрів, тягнуться до землі. Коли лідер наближається до поверхні, назустріч йому виростає короткий канал,
і~вздовж прокладеного лідером шляху проскакує зворотний (головний) розряд зі~струмом $10^{4}$--$10^{5}$~А. Канал діаметром кілька сантиметрів розігрівається до $\sim3\times10^{4}$~К;
стрімке розширення нагрітого газу й~дає грім. Плазма каналу живе лише соті--десяті частки секунди: потім вона охолоджується й~іони рекомбінують.

Усі ці об'єкти різняться на десятки порядків за концентрацією, але до всіх застосовні однакові поняття:
дебаївське екранування (§\,\ref{sec:debye-screening}), квазінейтральність (§\,\ref{sec:quasineutrality}),
рух у~магнітному полі (розд.~\ref{sec:particle-motion}). Докладніше про космічну плазму див.~\cite{gurnett2017introduction}.

\section{Лабораторна й~технічна плазма}\label{sec:nature-lab}

Лабораторну плазму отримують штучною іонізацією газу (способи перелічено в~§\,\ref{sec:what-is-plasma}), найчастіше електричним розрядом. Вона, як правило, холодніша й~густіша за космічну (виняток~--- гаряча плазма для синтезу), зате її можна вимірювати й~керувати нею
(рис.~\ref{fig:nature-lab}).

\begin{figure}[tb]\centering
	\begin{subfigure}[t]{0.45\linewidth}\centering
		\includegraphics[width=\linewidth]{discharge-tubes}
		\subcaption{Розрядні трубки з~He, Ne, Ar, Kr, Xe (Alchemist-hp, CC BY-SA 2.0 DE)}
	\end{subfigure}\hfill
	\begin{subfigure}[t]{0.45\linewidth}\centering
		\includegraphics[width=\linewidth]{plasma-globe}
		\subcaption{Плазмова куля (Supportstorm)}
	\end{subfigure}

	\medskip
	\begin{subfigure}[t]{0.45\linewidth}\centering
		\includegraphics[width=\linewidth]{welding-arc}
		\subcaption{Дугове зварювання (W.~M.~Plate~Jr.)}
	\end{subfigure}\hfill
	\begin{subfigure}[t]{0.45\linewidth}\centering
		\includegraphics[width=\linewidth]{tokamak-mast}
		\subcaption{Токамак MAST, Culham (Eye Steel Film, CC BY 2.0)}
	\end{subfigure}
	\caption{Лабораторна й~технічна плазма: газовий розряд~(а), плазмова куля~(б), електрична дуга~(в), токамак~(г).
	Знімки (б) і~(в)~--- суспільне надбання}
	\label{fig:nature-lab}
\end{figure}

\subsection*{Газові розряди}

У~\emphx{розрядній трубці} (рис.~\ref{fig:nature-lab}\,а) світиться слабоіонізований газ низького тиску: електричне поле розганяє вільні електрони, ті іонізують і~збуджують атоми,
а~колір дає лінійчатий спектр збуджених атомів (неон світить червоним, аргон~--- фіолетовим). Електрони набирають від поля енергію, а~іони, нерухливі й~важкі, її майже не отримують:
тому $T_e\sim1$~еВ, а~$T_i$~--- близько кімнатної, а~концентрація електронів у~тліючому розряді лише $\sim10^{10}~\mathrm{см}^{-3}$ (§\,\ref{sec:two-temperatures}). У~\emphx{плазмовій кулі}
(рис.~\ref{fig:nature-lab}\,б) високочастотна напруга на центральному електроді запалює розряд у~розрідженому газі, а~тонкі світні канали~--- це \emph{стримери}.
\emphx{Зварювальна дуга} (рис.~\ref{fig:nature-lab}\,в) навпаки, горить за високого тиску. Тут частих зіткнень вистачає, щоб вирівняти $T_e$ й~$T_i$ (порядку 1~еВ), тож плазма \emph{ізотермічна}, а~її
концентрація на кілька порядків вища, ніж у~трубці. Теорія розрядів викладена в~розд.~\ref{sec:discharge}.

Якщо потрібна чиста плазма без забруднення матеріалом електродів, користуються \emph{безелектродними} розрядами: високочастотним чи надвисокочастотним.
Продуваючи газ крізь розряд, отримують струмінь плазми; такий пристрій називають \emphx{плазмотроном}\index{Плазмотрон|textbf}. Температура струменя сягає $(1$--$3)\times10^4$~К, а~потужність пристрою~--- від десятків ватів до десятків мегаватів.
Для досліджень, де хвилі в~плазмі мають бути слабкими, будують спокійну плазму в~\emph{Q-машині} (від англ.\ \emph{quiescent}~--- спокійний): атоми цезію, потрапляючи на розжарену вольфрамову пластину, втрачають електрон і~стають
іонами, а~сама пластина випромінює електрони. Виходить майже повністю іонізована плазма з~температурою $\sim0{,}2$~еВ.

\subsection*{Плазма в~технологіях}

Плазма давно вийшла з~лабораторії (рис.~\ref{fig:nature-tech}). Почнімо з~теплових застосувань. У~\emphx{плазмовому різаку} (різновиді плазмотрона; рис.~\ref{fig:nature-tech}\,а)
електрична дуга стискається вузьким соплом і~потоком газу перетворюється на струмінь із~температурою $\sim2\times10^{4}$~К, що розплавляє метал і~здуває розплав;
так розрізають сталеві листи завтовшки до кількох сантиметрів. Той самий принцип, але зі~струменем, у~який вдувають порошок, використовують для \emph{плазмового напилення}
(рис.~\ref{fig:nature-tech}\,б): частинки металу чи кераміки розплавляються в~плазмі, вдаряються об~деталь і~утворюють покриття, скажімо теплозахисне для двигунів.

Інший напрям~--- \emph{хімічний}, а~не тепловий. У~газі низького тиску ($\sim1$--$100$~мТор) високочастотний розряд має $T_e$ кілька електронвольтів і~$n_e\sim10^{9}$--$10^{11}~\mathrm{см}^{-3}$~\cite{lieberman2005principles}, тоді як сам газ лишається майже холодним. Електрони розкладають молекули робочого газу (фтор- чи хлоровмісні) на хімічно активні радикали,
а~електричне поле біля поверхні розганяє іони перпендикулярно до пластини. Радикали реагують із~матеріалом, іони допомагають реакції й~задають їй напрям,
тож у~кремнії виходять вузькі глибокі канавки з~вертикальними стінками~--- так виготовляють мікросхеми. Таким \emphx{травленням}\index{Травлення плазмове|textbf} обробляють і~метали
(алюміній травлять хлоровмісною плазмою), і~скло; плазмою також очищають поверхні й~наносять тонкі плівки. Рис.~\ref{fig:nature-tech}\,в показує поверхню кремнію після плазмового травлення.

\begin{figure}[tb]\centering
	\begin{subfigure}[t]{0.45\linewidth}\centering
		\includegraphics[width=\linewidth]{plasma-cutting}
		\subcaption{Плазмове різання металу (D.~Baron Carral, CC0)}
	\end{subfigure}\hfill
	\begin{subfigure}[t]{0.45\linewidth}\centering
		\includegraphics[width=\linewidth]{plasma-spray}
		\subcaption{Плазмове напилення покриття на деталь двигуна (Optimi3e, CC BY-SA 4.0)}
	\end{subfigure}

	\medskip
	\begin{subfigure}[t]{0.45\linewidth}\centering
		\includegraphics[width=\linewidth]{plasma-etching}
		\subcaption{Поверхня кремнію після плазмового травлення, електронний мікроскоп (PolinLinB, CC BY 4.0)}
	\end{subfigure}\hfill
	\begin{subfigure}[t]{0.45\linewidth}\centering
		\includegraphics[width=\linewidth]{hall-thruster}
		\subcaption{Лабораторний двигун Холла на ксеноні, 6~кВт (NASA/JPL)}
	\end{subfigure}
	\caption{Плазма в~технологіях: різання~(а), напилення~(б), травлення~(в), космічний двигун~(г).
	Знімок (г)~--- суспільне надбання}
	\label{fig:nature-tech}
\end{figure}

Нарешті, \emph{космічні двигуни}. У~двигуні Холла (рис.~\ref{fig:nature-tech}\,г) іони ксенону розганяються електричним полем у~кільцевому каналі,
де радіальне магнітне поле утримує електрони, і~вони дрейфують у~схрещених полях $\vect{E}\times\vect{B}$ (розд.~\ref{sec:particle-motion}). Швидкість витікання сягає $10$--$30$~км/с~---
у~кілька разів більше, ніж у~хімічних ракетних двигунів, тож витрата палива для тієї самої зміни швидкості мала. Це далеко не весь перелік: газорозрядні лампи й~вивіски, дисплеї,
плазма олова, розігріта лазером, як джерело ультрафіолетового випромінювання для літографії, холодна плазма атмосферного тиску, яку досліджують для знезараження поверхонь і~води,
в~медицині та сільському господарстві. Спільне для всіх цих технологій: плазма дає змогу створити хімічно активні частинки без високої температури всього газу
або, навпаки, досягти температур, недосяжних для полум'я.

\subsection*{Гаряча плазма для термоядерного синтезу}

Окремий клас~--- гаряча плазма з~температурою порядку $10^{8}$~К, потрібна для реакцій термоядерного синтезу
(розд.~\ref{sec:fusion}). Контакт зі~стінкою миттєво охолодив би плазму (і~зруйнував би стінку), тому її \emph{утримують} магнітним полем (розд.~\ref{sec:magnetic-confinement}):
у~токамаку (рис.~\ref{fig:nature-lab}\,г), стелараторі чи пастці з~дзеркалами. У~видимому світлі світиться насамперед холодніша периферія плазми,
а~гаряче ядро ($T\sim10^{4}$~еВ) майже не світить: воно випромінює переважно в~інших діапазонах. Інший підхід~--- \emph{інерційний}: мішень стискають лазерними імпульсами так швидко,
що вона не встигає розлетітися. У~грудні 2022~р. в~Національному комплексі запалювання (NIF, США) вихід енергії синтезу вперше перевищив енергію лазерних імпульсів, що надійшли на мішень~\cite{llnl2025ignition},
а~у~токамаку JET 2023~р. у~D--T-розряді отримали рекордні 69~МДж за $\sim5$~с~\cite{eurofusion2024jet}. Будують і~міжнародний реактор ITER~\cite{iter2024baseline}.

\section{Схожі на вигляд, різні за природою}\label{sec:nature-pairs}

Інтуїція про плазму легко підводить: схожі на вигляд об'єкти бувають фізично різними. Розгляньмо три пари, на яких видно,
що саме треба вимірювати, щоб відрізнити одне від іншого.

\subsection*{Вогонь і~плазма}
Полум'я (рис.~\ref{fig:nature-flame}) гаряче, світиться й~містить заряджені частинки, тож його нерідко називають плазмою.

\begin{wrapstuff}[r,type=figure,width=0.28\linewidth,top=0]
	\includegraphics[width=\linewidth]{kerosene-flame}
	\caption{Полум'я гасової лампи (Arivumathi, CC0)}
	\label{fig:nature-flame}
\end{wrapstuff}
Це ламінарне полум'я: газ тече спокійно, без вихорів, тому видно його шарувату будову~--- блакитну зону внизу, де паливо реагує з~повітрям, і~вище яскраву жовту зону розжарених частинок сажі.
Світло дають саме вони й~збуджені радикали (CH, C$_2$), а~не плазма. Така будова задана хімією горіння, тобто реакціями палива з~киснем повітря. Заряджених частинок у~полум'ї дуже мало: лише кілька на мільярд нейтральних атомів і~молекул, і~утворюються вони хімічно, а~не нагріванням.
Чи достатньо їх, щоб газ поводився як плазма, покажуть три умови, виведені в~розд.~\ref{sec:plasma-basics}: ми перевіримо їх для полум'я
в~§\,\ref{sec:flame-check} і~побачимо, що воно лежить на межі.

\subsection*{Розрядна трубка й~дуга}

Неонова лампа й~зварювальна дуга обидві світяться й~обидві є газовим розрядом (рис.~\ref{fig:nature-lab}\,а,\,в), але фізично це різні плазми.
Різницю підсумовує табл.~\ref{tab:nature-tube-arc}: у~трубці низького тиску електрони рідко стикаються з~важкими частинками й~не віддають їм енергію, тому $T_e\gg T_i$;
у~дузі високого тиску зіткнень вистачає, щоб вирівняти температури.

\begin{table}[ht]
	\caption{Розряд низького тиску й~дуга: орієнтовні параметри}
	\label{tab:nature-tube-arc}
	\centering\small
	\begin{platbl}{colspec={X[1.4,l,m]X[c,m]X[c,m]}}
		& Розрядна трубка & Дуга \\
		Тиск & $\sim0{,}1$--$10$~Тор & $\gtrsim1$~атм \\
		$T_e$ & $\sim1$~еВ & $\sim1$~еВ \\
		$T_i$ & $\sim300$~К & $\sim1$~еВ \\
		$n_e$, $\mathrm{см}^{-3}$ & $\sim10^{10}$ & $10^{15}$--$10^{17}$ \\
		Ступінь іонізації & $\lesssim10^{-6}$ & відсотки й~більше \\
		Рівновага & $T_e\gg T_i$ & ізотермічна \\
	\end{platbl}
\end{table}

\subsection*{Сонце й~токамак}

І~те, й~те~--- гаряча плазма, у~якій відбувається (чи має відбуватися) термоядерний синтез, але плазму утримують принципово по-різному. Сонце тримає власна гравітація, а~в~токамаку~--- магнітне поле
(розд.~\ref{sec:magnetic-confinement}). Це визначає й~те, наскільки \enquote{щільно} виділяється енергія.

\begin{Estimate}[Скільки потужності дає кубометр Сонця]
Світність Сонця $L\approx3{,}8\times10^{26}$~Вт виділяється майже вся в~ядрі радіусом $0{,}2R_\odot\approx1{,}4\times10^{8}$~м, тобто в~об'ємі $V\approx1{,}1\times10^{25}~\mathrm{м}^3$.
Середня питома потужність $L/V\approx30$~Вт/м$^3$, як у~купі компосту. Для термоядерного реактора потрібно близько мегавата на кубометр плазми, тобто в~десятки тисяч разів більше.
\end{Estimate}
Сонце виграє розмірами й~часом: його ядро велике, а~тримається гравітацією мільярди років, до того ж протон-протонна реакція в~ньому дуже повільна. Реактору ж потрібна швидка реакція, а~суміш дейтерію з~тритієм реагує значно швидше, проте лише за вищої температури. Тому на Землі плазма мусить бути значно гарячішою
(близько $10^{8}$~К проти $1{,}5\times10^{7}$~К у~ядрі Сонця) і~горіти на іншому паливі (розд.~\ref{sec:fusion}).

\section{Що об'єднує всі ці об'єкти}\label{sec:nature-common}

Підсумуємо розглянуте таблицею й~діаграмою: плазми різняться за концентрацією $n$ й~температурою $T$ на багато порядків (табл.~\ref{tab:plasma-types}).

\begin{table}[ht]
	\caption{Орієнтовні параметри різних плазм (значення наведено за порядком величини).}
	\label{tab:plasma-types}
	\centering\small
	\begin{platbl}{colspec={X[2.4,l,m]X[c,m]X[c,m]}}
		Плазма & $n_e$, $\mathrm{см}^{-3}$ & $T$, еВ \\
		Іоносфера & $10^{5}$ & $0{,}1$ \\
		Міжзоряний газ & $10$ & $1$ \\
		Сонячна корона & $10^{8}$ & $10^{2}$ \\
		Ядро Сонця & $10^{26}$ & $10^{3}$ \\
		Токамак & $10^{14}$ & $10^{4}$ \\
		Сонячний вітер (1~а.о.) & $10$ & $10$ \\
		Газовий розряд & $10^{12}$ & $3$ \\
		Електрони в~металі & $10^{23}$ & вироджені, $E_F\sim8$ \\
	\end{platbl}
\end{table}

Концентрації в~таблиці відрізняються на 25 порядків (від $10$ до $10^{26}~\mathrm{см}^{-3}$), температури~--- на п'ять (від $0{,}1$ до $10^{4}$~еВ).
Електрони в~металі \emphx{вироджені}\index{Газ!вироджений|textbf}: їхня енергія визначається не теплом, а~принципом Паулі ($T\ll E_F$), де $E_F$~--- енергія Фермі, тобто найбільша енергія електронів при $T=0$;
такий газ потребує квантового опису (§\,\ref{sec:thomas-fermi}). Діаграму на рис.~\ref{fig:plasma-types} читають так: кожна точка~--- плазма з~даними $n$ і~$T$, а~лінії
відокремлюють області, де змінюється фізика: $N_D\sim1$~--- межа колективної поведінки (§\,\ref{sec:debye-sphere}), $T=E_F$~--- виродження, $T=m_ec^2$~--- релятивізм ($\hbar$~--- стала Планка, $c$~--- швидкість світла).
Величини $\lambda_{De}$, $N_D$ й~$\Gamma$ означимо в~розд.~\ref{sec:plasma-basics}.

\begin{figure}[ht]\centering
	\localinput{plasma_types.tikz}
	\caption{Діаграма $\lg n$--$\lg T$ видів плазми: лінії $N_D\sim1$ (орієнтовна межа слабкозв'язаної плазми; практично потрібно $N_D\gtrsim100$), $T=E_F$ (виродження)
	та $T=m_ec^2$ (релятивізм) обмежують область класичної ідеальної плазми $\max\{e^2n^{1/3},\hbar^2n^{2/3}/m_e\}\ll T\ll m_ec^2$}
	\label{fig:plasma-types}
\end{figure}

\begin{Simulation}
Побачити всі розглянуті об'єкти на одній діаграмі $\lg n$--$\lg T$, подивитися їхні параметри ($\lambda_{De}$, $N_D$, $\Gamma$, $\omega_{pe}$) і~задати власну пару $n$, $T$, щоб дізнатися, чи це ще колективна плазма, можна в~\simlink[Мапа плазм][інтерактивна діаграма $\lg n$--$\lg T$: об'єкти природи й~техніки з~розд.~0, лінії $N_D=1$, $\Gamma=1$, виродження й~релятивізму, калькулятор параметрів для довільних $n$ і~$T$]{nature_plasma_map.html}{інтерактивній мапі плазм}. Клацніть по знімку-об'єкту~--- побачите його фото, параметри й~режим; перетягніть \enquote{мій об'єкт} під лінію $N_D=1$~--- колективна поведінка зникне.
\end{Simulation}

Ми бачили плазму в~усіх масштабах: від каналу іскри завтовшки кілька сантиметрів до світлових років, від кімнатної температури іонів до сотень мільйонів кельвінів. Три властивості,
що повторюються в~більшості прикладів, і~є змістом нашого курсу.

\begin{itemize}
	\item \emph{Колективність і~квазінейтральність.} Плазма квазінейтральна на масштабах, більших за~$\lambda_D$, але заряди в~ній взаємодіють усереднено, на великих відстанях. Звідси екранування й~власні коливання
	(розд.~\ref{sec:plasma-basics}, \ref{sec:plasma-waves}).
	\item \emph{Магнітне поле.} Воно вказує частинкам траєкторію: звідси дрейфи, магнітні пастки й~вморожене поле (розд.~\ref{sec:particle-motion}, \ref{sec:mhd}, \ref{sec:magnetic-confinement}).
	\item \emph{Роль зіткнень.} У~розрідженій плазмі (космос, розряд низького тиску) зіткнення рідкісні, тож потрібні кінетичний опис і~теорія хвиль та нестійкостей; у~щільній (дуга, полум'я)
	зіткнення часті й~визначають перенос (розд.~\ref{sec:plasma-collisions}, \ref{sec:plasma-kinetic}, \ref{sec:instability}).
\end{itemize}

Решта книги~--- про те, як ці властивості виражаються формулами. Спершу, у~розд.~\ref{sec:plasma-basics}, з'ясуємо, що саме відрізняє плазму від газу; далі~--- хвилі в~холодній плазмі (розд.~\ref{sec:plasma-waves}),
рух частинок (розд.~\ref{sec:particle-motion}), зіткнення (розд.~\ref{sec:plasma-collisions}), гідродинамічний опис (розд.~\ref{sec:mhd}), хвилі в~замагніченій плазмі (розд.~\ref{sec:mhd-waves}),
кінетика (розд.~\ref{sec:plasma-kinetic}), нестійкості (розд.~\ref{sec:instability}), утримання плазми й~синтез (розд.~\ref{sec:magnetic-confinement}, \ref{sec:fusion}), газовий розряд (розд.~\ref{sec:discharge})
і~підсумкова карта масштабів (розд.~\ref{sec:plasma-scales}).

\Needspace{8\baselineskip}
\begin{Summary}
\begin{itemize}
\item Понад 99\,\% видимої речовини у~Всесвіті~--- плазма: від первинної воднево-гелієвої плазми раннього Всесвіту (рекомбінація за $T\approx0{,}3$~еВ) до зір, туманностей, міжгалактичного газу й~акреційних дисків чорних дір.
\item Природна плазма охоплює концентрації від $\sim10^{-5}$ до $10^{26}~\mathrm{см}^{-3}$ й~розміри від сантиметрів до світлових років; у~розрідженій космічній плазмі зіткнень мало, а~структури задає магнітне поле
(петлі корони, сяйво, вморожені волокна туманностей).
\item Лабораторна плазма: розряд низького тиску нерівноважний ($T_e\gg T_i$), дуга ізотермічна; плазмотрони, різання, напилення, травлення, двигуни Холла; гаряча плазма для синтезу, яку утримує магнітне поле.
\item Плазма~--- квазінейтральний іонізований газ із~колективною поведінкою; низько- й~високотемпературну плазму розрізняють за $T$ (межі $\sim10^5$ і~$\sim10^6$~К). Полум'я~--- гранична слабоіонізована плазма ($\alpha\sim10^{-9}$): зіткнення з~нейтралами домінують над плазмовими процесами (перевірка~--- §\,\ref{sec:flame-check}).
\item Об'єкти розрізняють за $n$, $T_e$ і~$T_i$, ступенем іонізації та способом утримання (гравітація, магнітне поле, інерція, стінки).
\end{itemize}
\end{Summary}

\begin{Questions}
\begin{quizlist}
\item Чому Всесвіт був непрозорим до рекомбінації і~став прозорим після неї?

\item Чому водень рекомбінує за $T\sim0{,}3$~еВ, хоча потенціал іонізації становить $13{,}6$~еВ?

\item Чому петлі в~короні Сонця повторюють силові лінії магнітного поля?

\item Чим плазма розрядної трубки відрізняється від плазми дуги за $T_e$, $T_i$ і~ступенем іонізації? Чому так?

\item Чому полум'я~--- слабкоіонізована плазма? Яка величина показує, що колективні процеси в~ньому придушені? (Підказка: §\,\ref{sec:flame-check}.)

\item Чому на Землі для термоядерного реактора потрібна вища температура плазми, ніж у~ядрі Сонця?

\item Чому іоносфера відбиває короткі радіохвилі, а~ультракороткі проходять крізь неї?

\item Чому в~міжзоряному газі температури електронів та іонів можуть помітно відрізнятися?

\item Чому світло, яке бачить Event Horizon Telescope, випромінює не сама чорна діра, а~плазма довкола неї?
\end{quizlist}
\end{Questions}
